%=============================================================================!
% 3.4  POTENTIAL ENERGY AND THE CONSERVATION OF ENERGY                        !
%=============================================================================!
\section{Potential energy and the conservation of energy}
\label{sec:en-potential}

A ball thrown upwards slows down, stops and falls back, arriving at the
hand with the speed it left with. Its kinetic energy has gone and come
back, as though it had been stored somewhere on the way up. This section
makes the store precise for a position-dependent force along a line, and
shows that Newton's second law then keeps the total fixed.

\subsection*{Potential energy}

The weight of a body and the pull of a spring both depend only on where
the body is, not on how fast it moves or on the time. For such a force
$F(x)$ along a line, \cref{eq:en-work-line} gives the work from $x_A$ to
$x_B$ as an integral fixed by the two ends. Choose a fixed reference
point $x_0$, which need not be where the motion starts, and define the
potential energy
\begin{equation}
   U(x) = -\int_{x_0}^{x} F(x')\,\dd x' ,
   \label{eq:en-potential}
\end{equation}
minus the work the force does as its body moves from $x_0$ to $x$. Two
rules for integrals follow from writing each one, as in
\cref{sec:en-varying}, as the change $G(b) - G(a)$ in a function with $G'
= F$: reversing the limits changes the sign, since $G(a) - G(b) =
-[G(b) - G(a)]$, and an integral can be split at any point $c$, inside the
range or not, since $[G(c) - G(a)] + [G(b) - G(c)] = G(b) - G(a)$. Then,
splitting the integral at $x_0$ and reversing the first part,
\begin{align*}
   W_{A\to B} &= \int_{x_A}^{x_B} F\,\dd x
   = \int_{x_A}^{x_0} F\,\dd x + \int_{x_0}^{x_B} F\,\dd x
      && \text{going by way of $x_0$}\\
   &= -\int_{x_0}^{x_A} F\,\dd x + \int_{x_0}^{x_B} F\,\dd x
      && \text{reversing the first integral}\\
   &= U(x_A) - U(x_B)
      && \text{by \cref{eq:en-potential}.}
\end{align*}
The work done by the force is the drop in potential energy. The force
uses up potential energy when it does positive work, and stores it when
it does negative work, as the weight does while the ball rises.

Moving the reference point from $x_0$ to $x_1$ adds the same constant to
$U$ everywhere, since splitting the integral at $x_0$ gives
$-\int_{x_1}^{x} F\,\dd x' = -\int_{x_1}^{x_0} F\,\dd x' + U(x)$, whose
first term does not depend on $x$; this changes no drop in $U$. Only differences of potential
energy have a physical meaning, and we choose $x_0$ wherever it is
convenient. The force follows from $U$: by the fundamental theorem of
calculus, read as in \cref{sec:en-varying}, the derivative of
$\int_{x_0}^{x}F(x')\,\dd x'$ with respect to its upper limit is $F(x)$,
so
\begin{equation}
   F(x) = -\frac{\dd U}{\dd x} .
   \label{eq:en-force-1d}
\end{equation}
The force points the way $U$ decreases, and the steeper the graph of $U$,
the stronger the force.

Two potential energies follow at once. For a body near the ground, with
$x$ its height measured upwards, the weight is $F = -mg$, and with $x_0 =
0$,
\begin{equation*}
   U(x) = -\int_0^x(-mg)\,\dd x' = mgx .
\end{equation*}
For a spring, $F = -kx$, and with $x_0 = 0$ at the natural length,
\begin{equation*}
   U(x) = -\int_0^x(-kx')\,\dd x' = \tfrac12 k x^2 ,
\end{equation*}
the work of \cref{sec:en-varying} now stored in the stretched spring.

\subsection*{Conservation of energy}

When the net force is of this kind, the work-energy theorem $K_B - K_A = W_{A\to
B}$ and the drop $W_{A\to B} = U_A - U_B$ combine into $K_B - K_A = U_A -
U_B$, or, moving $U_B$ to the left and $K_A$ to the right,
\begin{equation*}
   K_B + U_B = K_A + U_A .
\end{equation*}
The sum of the two energies is the same at every pair of points along the
motion. We call it the total energy, $E = K + U$. The argument can also
be run as a rate, which shows exactly where Newton's law enters.

\begin{proposition}
\label{prop:en-energy-1d}
Let a body of mass $m$ move along a line under a net force $F(x) = -\dd
U/\dd x$ that depends only on its position. Then its total energy
\begin{equation*}
   E = \tfrac12 m v^2 + U(x)
\end{equation*}
keeps the same value throughout the motion.
\end{proposition}

\begin{proof}
Differentiate $E$ along the motion, using the chain rule on each term:
$\frac12 mv^2$ changes at the rate $mv\,\dd v/\dd t$, and $U(x(t))$ at the
rate $(\dd U/\dd x)\,(\dd x/\dd t) = -F v$. Hence
\begin{align*}
   \frac{\dd E}{\dd t} &= m v\,\frac{\dd v}{\dd t} - F v
   = v\,\Bigl(m\,\frac{\dd v}{\dd t} - F\Bigr) = 0 ,
\end{align*}
because the bracket vanishes by the second law, $m\,\dd v/\dd t = F$. A
quantity whose rate of change is zero does not change.
\end{proof}

Newton's second law enters at a single point: it turns the rate of change
of the kinetic energy into the power of the force, which the potential
energy then exactly cancels. Energy is therefore conserved, in the sense
of \cref{sec:ne-bodies}, and \cref{prop:en-energy-1d} is the conservation
of energy. Forces that have a potential energy, and so
conserve $E$, are called conservative.

\subsection*{Two checks}

Energy reproduces the results of the earlier chapters with no equation
of motion solved. A ball of \SI{0.5}{\kilogram} thrown up at
\SI{12}{\metre\per\second} starts at $x = 0$, where $U = 0$, with $E = K =
\frac12\times0.5\times12^2 = \SI{36}{\joule}$. At the top of its flight
$K = 0$, so all of $E$ is potential energy, $mgx_{\max} = E$, and
\begin{equation*}
   x_{\max} = \frac{E}{mg} = \frac{36}{0.5\times9.80665}
   = \SI{7.342}{\metre} ,
\end{equation*}
the height found in \cref{sec:mo-integrating}. Halfway up, at
\SI{3.671}{\metre}, $U = mgx$ is half its value at the top, so half of $E$
is potential energy and half kinetic: $\frac12 mv^2 = \SI{18}{\joule}$,
half the kinetic energy of the throw, so $v^2$ is half of $12^2$ and $v =
12/\sqrt2 = \SI{8.485}{\metre\per\second}$. \Cref{fig:en-ball} follows $K$ and $U$
through the flight: they trade with each other while their sum stays at
\SI{36}{\joule}. Plotted against height, $U = mgx$ is a straight line and
so is $K = E - mgx$, which reaches zero exactly at the top.

For the block on a spring, at a turning point $x = \pm A$ the block is at
rest and $E = \frac12 kA^2$; at the centre $U = 0$ is at its smallest, so
$K = E - U$ and the speed are at their largest, and $E = \frac12
mv_{\max}^2$. Equating the two,
\begin{equation*}
   \tfrac12 m v_{\max}^2 = \tfrac12 k A^2
   \quad\Longrightarrow\quad
   v_{\max} = A\sqrt{\frac{k}{m}} = A\omega ,
\end{equation*}
with $\omega$ from \cref{eq:ne-spring-omega}, as in \cref{sec:ne-spring}.

\begin{figure}[htb]
\centering
\includegraphics{ch03_energy/ball}
\caption{The energies of a ball of \SI{0.5}{\kilogram} thrown straight up
at \SI{12}{\metre\per\second}. (a) Kinetic energy $K$, potential energy
$U = mgx$ and their sum $E$ against time; the dotted line marks the top.
(b) The same against height: $K$ and $U$ are straight lines, and $K$
reaches zero at the top, \SI{7.342}{\metre}.}
\label{fig:en-ball}
\end{figure}

\begin{keyidea}
A force that depends only on position along a line has a potential energy
$U(x) = -\int_{x_0}^{x} F(x')\,\dd x'$, with $F = -\dd U/\dd x$. Newton's second law then
keeps the total energy $E = K + U$ constant: kinetic and potential energy
are traded, never their sum.
\end{keyidea}

\notebook{03}{follow $K$, $U$ and $E$ through the throw with a time
cursor}

\begin{selfcheck}
A ball is thrown at \SI{12}{\metre\per\second} from a balcony
\SI{5}{\metre} above the ground. With what speed does it reach the
ground? Does the answer depend on whether it is thrown upwards, sideways
or downwards?
\end{selfcheck}
